\section{Literature Review}

\subsection{Chart Understanding with Vision-Language Models}

A growing body of work has explored how Vision-Language Models can be
applied to chart understanding. ChartQA \cite{masry2022chartqa}
introduces a benchmark where models are required to answer questions
about charts, combining visual perception with logical reasoning.
DePlot \cite{liu2023deplot} takes a different approach by converting
chart images into structured tables, which can then be used for
numerical reasoning. More recently, UniChart \cite{masry2023unichart}
proposes a unified framework for chart-related tasks through
pretraining. Together, these studies suggest that VLMs are effective at
extracting information from charts and performing reasoning tasks.
However, the focus of this line of work is on understanding what a
chart shows, rather than evaluating whether the chart itself may be
misleading. Detecting misleading patterns requires a fundamentally
different perspective, one that considers how visual encoding choices
may distort the underlying data rather than simply interpreting what is
presented. Our work addresses this gap by treating chart analysis as an
auditing task rather than a comprehension task.

\subsection{Misleading Visualization Detection}

Misleading visualizations have traditionally been studied from the
perspective of human perception and design principles.
\citet{pandey2015deceptive} demonstrate that truncated axes can
significantly affect how differences are perceived, while
\citet{lisnic2023misleading} provide a broader taxonomy of misleading
techniques observed in real-world media. The Misviz benchmark
\cite{tonglet2025misviz} represents the first large-scale effort to
formalize this problem for machine learning, providing charts annotated
across 12 misleader types. However, existing evaluations have focused
primarily on binary detection, leaving per-type detection ability
largely unexplored. No prior automated system has been shown to
reliably identify specific misleader types, particularly in realistic
financial contexts. Our work directly addresses this gap by evaluating
and improving per-type detection performance as the primary metric.

\subsection{Tool-Augmented Language Models}

Tool-augmented language models have shown strong performance across a
range of reasoning tasks. Toolformer \cite{schick2023toolformer}
demonstrates that models can learn to invoke external tools to improve
task performance, while ReAct \cite{yao2023react} combines step-by-step
reasoning with tool usage in an interleaved process. ViperGPT
\cite{suris2023vipergpt} extends this idea to visual reasoning by
enabling models to execute programmatically generated code. While these
approaches have proven effective in text-based and programmatic
settings, their applicability to chart visual reasoning is not well
understood. Adding external information such as OCR outputs to a visual
prompt may not always help and can sometimes interfere with the model's
visual attention. Whether and how tool augmentation should be integrated
into chart analysis pipelines remains an open question that our work
directly investigates.

\subsection{SEC Financial Disclosure Analysis}

In the financial domain, prior research has focused predominantly on
textual and numerical analysis. FinQA \cite{chen2021finqa} introduces a
dataset for numerical reasoning over financial documents, and
\citet{black2021nongaap} study the consistency and comparability of
non-GAAP earnings disclosures, highlighting ongoing regulatory
concerns. However, these studies largely overlook visual elements such
as charts, despite their widespread use in financial filings. No prior
work has examined chart-level compliance with SEC regulations such as
Regulation~G or Item~10(e) of Regulation~S-K in an automated way. Our
work addresses this gap by extending VLM-based chart analysis to real
SEC 10-K filings and validating findings against actual regulatory
comment letters.

% ============================================================
